\scanpage{8}
\deleted{
\uline{Cor} $\mathcal A_1$ est stable par primitivation canonique (i.e. primitive d'augmentation nulle) contient $\log(1+t_1)$ (primitive de $\frac1{1+t_1}$) donc $\exp.t_1$ (série inverse), les fonctions $\cos t_1$, $\sin t_1$, $\operatorname{tg}t_1$ sont $\in\mathcal A_1$ $\ldots$

\uline{Lemme}\quad\deleted{\unsure{g) Soit}} Soit $X$ un\inserted{$\mathcal A$-} schéma formel (lisse)\inserted{de dim $n=p+q$}, $\Phi$ un ``système de Pfaff de degré $p$ complètement intégrable sur $X$'', \deleted{\unsure{$\Phi$}}
\deleted{\unsure{et un quotient}\quad \unsure{Soient} $Y,Z$ des schémas formels de dim $p$ resp $q$,
$f:X\to Y$ et $i:Z\to X$\inserted{des $\mathcal A$-morphismes}, tels que $T(p)(0)$ soit un \deleted{\unsure{isom}} $\Phi(0)\unsure{-|\to}T(Y)_0$ et que \deleted{\unsure{$T(i)(0)$}} le composé \deleted{\unsure{$T(X)_0$}}
\[
T(Z)_{\deleted{\unsure{(0)}}0}\overset{T(i)_0}{\unsure{-|\longrightarrow}}T(X)_{\deleted{\unsure{(0)}}0}\overset{\mathrm{can}}{\longrightarrow}T(X)_0/\Phi(0)
\]
soit un isom. On sait qu'il existe alors un isomorphisme\inserted{formel ($\varphi$ est néc. $\mathcal A$-isom)} unique
\[
X\overset{\varphi}{\simeq}Y\times Z
\]
\unsure{transformant $\Phi$ en produit} tel que $\operatorname{pr}_1\circ\varphi=p$, $\varphi\circ i=(z\mapsto(0,z))$)}

Soit $Z$\inserted{$\mathcal A$-} schéma formel\inserted{(lisse)} de dim $q$, $i:Z\unsure{-|\to}X$ une $\mathcal A$-immersion telle que le composé
\[
T(\unsure{Z|V})_0\overset{T(i)_0}{\unsure{-|\longrightarrow}}T(X)_0\longrightarrow T(X)_0/\Phi(0)
\]
soit un isom. On sait qu'il existe alors un unique morphisme formel
\[
X\xrightarrow{\pi}Z
\]
tel que $\pi\circ i=\operatorname{id}$, et que ``$\pi$ soit \deleted{\unsure{constant sur}}\inserted{compatible avec} le feuilletage défini par $\Phi$''. Conditions équivalentes

(i) $\pi$ est un $\mathcal A$-morphisme

(ii) $\exists$ $\mathcal A$-isom $X\simeq B\times Y$, tel que le feuilletage formel de $X$ défini par $\Phi$ soit le feuilletage défini par la projection $\operatorname{pr}_1:X\to B$. On dit
}
\figure{\begin{tikzpicture}[x=1cm,y=1cm]
\draw (.2,8)--(2.2,4.8)--(4.8,.3);\draw (2.4,7.9)--(5.3,4.6);\draw (4.6,7.7)--(5.7,5.2);\draw (0,6.2)--(0,4.1)..controls(0,4)and(.2,4)..(.3,4);\draw (0,6.5)--(5.8,6.45);\draw (.1,4.4)--(4.8,4.55);
\end{tikzpicture}}
\margin{7}
